% !TeX program = xelatex
% 编译：xelatex 03-线性时不变(LTI)系统的定义.tex（VS Code 的 LaTeX Workshop 会读上一行，自动用 xelatex）
% 来源：2022 浙江大学《信号与系统》胡浩基老师，第二章「LTI 系统时域分析」开篇，教材第 35 页，视频 00:00--21:22（整段视频；视频在老师讲到「LTI 很简单」的核心句时结束）。
\documentclass[UTF8,a4paper,11pt]{ctexart}
\usepackage{amsmath,amssymb}
\usepackage{geometry}
\geometry{margin=2.4cm}
\usepackage{xcolor}
\usepackage{mdframed}
\usepackage{booktabs}
\usepackage{needspace}
\usepackage{array}
\usepackage{longtable}
\usepackage{enumitem}
\usepackage{tikz}
\usetikzlibrary{arrows.meta,positioning,calc}
\usepackage{float}
\usepackage{fancyhdr}
% 页眉页脚清空，只在右下角保留页号（含首页）
\fancypagestyle{plain}{%
  \fancyhf{}%
  \fancyfoot[R]{\thepage}%
  \renewcommand{\headrulewidth}{0pt}%
  \renewcommand{\footrulewidth}{0pt}}
\usepackage[hidelinks,unicode,bookmarksnumbered,bookmarksopen=true,bookmarksopenlevel=2,pdftitle={线性时不变（LTI）系统的定义}]{hyperref}
\usepackage{bookmark}% XeTeX 下 hyperref 自带的大纲不生效，用 bookmark 宏包生成 PDF 书签（需编译两遍）

\setCJKmainfont{Songti SC}[BoldFont={Songti SC Bold}]
\setCJKsansfont{PingFang SC}

\newmdenv[backgroundcolor=blue!4,linecolor=blue!50!black,linewidth=1pt,innertopmargin=6pt,innerbottommargin=6pt]{rememberbox}
\newenvironment{remember}{\begin{rememberbox}\textbf{重点}\par\smallskip}{\end{rememberbox}}
\newmdenv[backgroundcolor=yellow!8,linecolor=orange!70!black,linewidth=1pt,innertopmargin=6pt,innerbottommargin=6pt,nobreak=true]{notebox}
\newenvironment{note}{\begin{notebox}\textbf{说明}\par\smallskip}{\end{notebox}}
\newmdenv[backgroundcolor=gray!8,linecolor=gray!60,linewidth=0.8pt,innertopmargin=5pt,innerbottommargin=5pt]{skipbox}
\newenvironment{skippable}{\begin{skipbox}\textbf{可跳过}\par\smallskip}{\end{skipbox}}

% ---------------------------------------------------------------- 画图用的小宏
% 一行「输入 -> 系统 -> 输出」。#1 纵坐标，#2 输入，#3 输出，#4 右侧备注（灰字）
\newcommand{\sysline}[4]{%
  \node[anchor=east] at (2.9,#1) {#2};
  \draw[->,thick] (3.0,#1) -- (3.8,#1);
  \node[draw,rounded corners=2pt,fill=blue!6,minimum width=1.3cm,minimum height=0.62cm] at (4.45,#1) {\small 系统};
  \draw[->,thick] (5.1,#1) -- (5.9,#1);
  \node[anchor=west] at (6.0,#1) {#3};
  \node[anchor=west,gray!55!black,font=\footnotesize] at (9.4,#1) {#4};}

% ---------------------------------------------------------------- 标题层级
% 章(section)、小节(subsection)：深蓝；三级标题(subsubsection)、小标题(paragraph)：青色；例题：橙色徽章。
\definecolor{secblue}{RGB}{31,78,121}
\definecolor{exorange}{RGB}{190,90,10}
\definecolor{ptteal}{RGB}{0,110,110}
\ctexset{
  section={format=\raggedright\Large\bfseries\color{secblue},
           beforeskip=3.2ex plus 1ex minus .5ex, afterskip=1.4ex,
           aftertitle={\par\nobreak\vspace{-1.1ex}\noindent\textcolor{secblue}{\rule{\linewidth}{0.8pt}}\par\nobreak\vspace{0.6ex}}},
  subsection={format=\raggedright\large\bfseries\color{secblue},
              beforeskip=2.4ex plus .8ex minus .4ex, afterskip=1.0ex},
  subsubsection={format=\raggedright\normalsize\bfseries\color{ptteal},
                 beforeskip=1.8ex plus .5ex minus .3ex, afterskip=0.8ex}
}
% 例题标题：\example{标题}{结论}{}（第三个参数是标题后的灰色小字附注，一般留空）。左边橙色实心「例 n」，右边结论标签；同时进 PDF 大纲（第 3 级）。
\newcounter{ex}[section]
\newcommand{\example}[3]{%
  \refstepcounter{ex}\phantomsection
  \addcontentsline{toc}{subsubsection}{例 \theex\ #1（#2）}%
  \par\addvspace{1.6ex}\noindent
  {\setlength{\fboxsep}{3pt}%
   \colorbox{exorange}{\textcolor{white}{\textbf{例~\theex}}}\ %
   \textcolor{exorange}{\textbf{#1}}\ {\footnotesize\textcolor{gray!70!black}{#3}}%
   \hfill\colorbox{exorange!12}{\small\textcolor{exorange}{#2}}\par}%
  \nobreak\vskip0.6ex\relax\csname @afterheading\endcsname}
% 小标题（接着同一段往下写，不进大纲）：青色三角 + 粗体。把原来的黑色粗体 \paragraph 统一换成它，避免被当成章节。
\newcommand{\point}[1]{\par\addvspace{1ex}\noindent\textcolor{ptteal}{$\blacktriangleright$\,\textbf{#1}}\ }
\renewcommand{\paragraph}[1]{\point{#1}}

\title{线性时不变（LTI）系统的定义\\\large 浙江大学 2022《信号与系统》（胡浩基）第二章开篇（教材第 35 页）听课笔记}
\author{}
\date{}

\begin{document}
\pagestyle{plain}
\maketitle

%==================================================================
\section{LTI 是什么}
第二章「LTI 系统的时域分析」（教材第 35 页）的研究对象是 LTI 系统，先给出它的定义。

\subsection{名字}
\textbf{LTI} 是 Linear Time-Invariant 的首字母：L 取自 Linear（线性），T、I 取自 Time-Invariant（时不变）。后面再加上 System（系统），就是 \textbf{LTI 系统}。

\begin{figure}[H]
\centering
\begin{tikzpicture}[>=stealth]
  \node[draw,rounded corners=3pt,fill=blue!6,minimum width=3.2cm,minimum height=1.0cm,align=center] (a) at (0,0) {线性\\{\small \textbf{L}inear}};
  \node[draw,rounded corners=3pt,fill=blue!6,minimum width=3.8cm,minimum height=1.0cm,align=center] (b) at (4.4,0) {时不变\\{\small \textbf{T}ime-\textbf{I}nvariant}};
  \node[draw,rounded corners=3pt,fill=blue!6,minimum width=3.2cm,minimum height=1.0cm,align=center] (c) at (8.8,0) {系统\\{\small \textbf{S}ystem}};
  \node[font=\Large\bfseries] (l) at (3.4,-2.0) {L};
  \node[font=\Large\bfseries] (ti) at (4.2,-2.0) {T\,I};
  \node[font=\Large\bfseries] (s) at (5.2,-2.0) {系统};
  \draw[->,thick] (a.south) -- (l.north);
  \draw[->,thick] (b.south) -- (ti.north);
  \draw[->,thick] (c.south) -- (s.north);
  \node[anchor=west,gray!55!black,font=\footnotesize,align=left] at (6.4,-2.25) {中文「线性时不变系统」和\\英文 Linear Time-Invariant System\\都叫 LTI 系统};
\end{tikzpicture}
\caption{L 取自 Linear，T、I 取自 Time-Invariant；LTI 本身只是「线性时不变」，加上「系统」才是 LTI 系统。}
\end{figure}

\begin{remember}
\textbf{如果一个系统既线性、又时不变，就叫它线性时不变系统}（Linear Time-Invariant System，简称 LTI）。两个条件必须\textbf{同时}满足：只满足一个、或两个都不满足，都不是 LTI。
\end{remember}
线性指齐次性 $+$ 叠加性，时不变指输入平移多少、输出也平移多少（定义见《典型的系统》）。

%==================================================================
\needspace{11cm}
\section{六个性质的英文对照}
六个性质的中英文及反面：

\begin{center}
\small
\begin{tabular}{@{}lll@{}}
\toprule
性质（中文） & 英文 & 反面（英文）\\
\midrule
线性系统 & Linear System & non-linear\\
时不变系统 & Time-Invariant System & time-variant\\
因果系统 & Causal System & non-causal（也有写作 uncausal 的）\\
无记忆系统 & Memoryless System & memory（有记忆）\\
可逆系统 & Invertible System & non-invertible \\
稳定系统 & Stable System & unstable\\
\bottomrule
\end{tabular}
\end{center}

三处英文要注意：
\begin{itemize}[nosep]
\item 因果的反面：常见写法 \texttt{non-causal}，也有写作 \texttt{uncausal} 的。
\item 不可逆写作 \texttt{non-invertible}；时变也写作 \texttt{time-varying}。
\item 可逆的拼写是 \texttt{Invertible}（-ible 结尾），不是 \texttt{Invertable}。
\end{itemize}

%==================================================================
\section{为什么研究 LTI 系统}
LTI 系统是本课程的重点，研究它有三个理由：\textcircled{1} 为什么研究\textbf{线性}；\textcircled{2} 为什么研究\textbf{时不变}；\textcircled{3} LTI 本身\textbf{很简单}。前两个说明「这两个假设为什么合理」，第三个说明「它为什么好算」。

\subsection{理由一：线性}
\subsubsection{世界上没有绝对线性的系统}
以喇叭为例（输入电流、输出声音）。\textbf{齐次性做不到}：通 1、2、3 单位电流，声音可以分别大 1、2、3 倍；但通 100、10000 单位电流时，声音不可能相应放大 100、10000 倍——任何物理系统都有极限。\textbf{叠加性也做不到}：一个人说话喇叭能放出，两个人、三个人叠加，总有一次放不下。

\begin{figure}[H]
\centering
\begin{tikzpicture}[>=stealth,x=1.0cm,y=0.8cm]
  \draw[->] (0,0) -- (7.0,0) node[right] {输入电流};
  \draw[->] (0,0) -- (0,5.6) node[above] {声音大小};
  \foreach \k in {1,...,6} { \draw (\k,0.08) -- (\k,-0.08) node[below,font=\scriptsize,fill=white,inner sep=1pt] {$\k$}; }
  \foreach \y in {1,...,5} { \draw (0.08,\y) -- (-0.08,\y) node[left,font=\scriptsize,fill=white,inner sep=1pt] {$\y$}; }
  \node[below left,font=\scriptsize] at (0,0) {$0$};
  % 理想线性：输出始终与输入成正比
  \draw[thick,dashed,gray!70!black] (0,0) -- (5.2,5.2) node[right,font=\small] {理想线性：始终成正比};
  % 实际：小输入时与直线重合，大输入时饱和
  \draw[very thick,blue!60!black] (0,0) -- (3,3);
  \draw[very thick,blue!60!black] plot[smooth] coordinates {(3,3) (3.5,3.42) (4,3.7) (5,3.95) (6,4.05)};
  \node[blue!60!black,anchor=west,font=\small] at (6.05,4.05) {实际：放不大了};
  \draw[dotted,gray] (3,0) -- (3,3) -- (0,3);
\end{tikzpicture}
\caption{喇叭的输入--输出关系（示意，相对单位）。输入不超过 $3$ 时，实际曲线与直线重合，输入翻几倍、声音就大几倍；输入再大，声音不再跟着变大（饱和），齐次性被打破。}
\end{figure}

\begin{remember}
世界上没有绝对意义的线性系统：齐次性与叠加性在足够大的输入上都会被物理极限打破。
\end{remember}

\subsubsection{线性是合理的近似}
结论：\textbf{线性只是对现实问题的近似}。理由不是「世界上有很多线性系统」，而是：现实系统太复杂，\textbf{只能用最简单的模型——线性——去近似}。

\begin{remember}
研究线性系统的理由：没有绝对的线性系统，但线性模型是最简单的近似，工程上常用它分析复杂问题。
\end{remember}

\example{摩擦力公式~\texorpdfstring{$F=\mu N$}{F=μN}}{线性近似}{}
\textbf{模型}：压力 $N$ 与摩擦力 $F$ 的经验公式
\[
F=\mu N,
\]
它不是从基本规律推出的，而是经验近似。

\paragraph{推导。}设 $F=S(N)$，在 $N=0$ 处作\textbf{泰勒展开}(Taylor expansion)：
\[
F=S(N)\approx S(0)+S'(0)\,N+\frac{S''(0)}{2!}N^2+\cdots
\]
假设没有压力就没有摩擦力（$S(0)=0$），且二次以上的项都近似为 $0$，则
\[
F\approx S'(0)\,N.
\]
（泰勒展开：光滑函数在一点附近可写成「直线 + 弯 + 更弯」的叠加；线性近似只留第一项。）记 $\mu=S'(0)$：\textbf{$\mu$ 就是曲线在 $N=0$ 处的斜率}。

\paragraph{适用范围。}$\mu$ 并非常数：静摩擦时要使劲拉，$\mu$ 较大；高速运动时 $\mu$ 又变大。\textbf{只有运动在合理范围内，$\mu$ 才近似为常数}。真实的 $F$--$N$ 关系是一条弯曲的曲线：小 $N$ 陡、中间平缓、大 $N$ 又陡；只在中间取一小段，才能近似成直线（见下图）。

\begin{figure}[H]
\centering
\begin{tikzpicture}[>=stealth,x=0.92cm,y=0.78cm]
  \draw[->] (0,0) -- (5.2,0) node[below] {$N$（压力）};
  \draw[->] (0,0) -- (0,4.4) node[left] {$F$（摩擦力）};
  % 真实的曲线：从原点几乎竖直上冲，拐弯后是平缓的长坡，过了右边的点又往上翘
  \draw[very thick,blue!60!black] plot[smooth] coordinates {(0,0) (0.1,0.8) (0.3,1.35) (0.7,1.6) (1.6,1.8) (2.6,2.05) (3.3,2.5) (4.0,3.9)};
  % 中间那一点的切线，斜率记作 μ
  \draw[thick,orange!85!black] (0.6,1.58) -- (3.4,2.22) node[right] {$\mu$};
  \fill[orange!85!black] (1.6,1.8) circle (1.8pt);
  \fill[blue!60!black] (3.3,2.5) circle (1.8pt);
  \node[orange!85!black,anchor=west,font=\small] at (2.2,1.45) {只在这一小段近似成直线，斜率就是 $\mu$};
  \node[blue!60!black,anchor=west,font=\small] at (2.2,0.85) {真实关系：弯曲的曲线};
  \node[gray!55!black,font=\footnotesize,anchor=west] at (0.15,3.6) {$N$ 小：陡（要使劲拉）};
  \node[gray!55!black,font=\footnotesize,anchor=west] at (3.4,4.3) {$N$ 很大：又陡};
\end{tikzpicture}
\caption{摩擦力和压力的关系（形状示意）：从原点几乎竖直上冲，拐弯后平缓，过了右边的点又往上翘。\textbf{只在中间一段，$F$ 才近似正比于 $N$。}}
\end{figure}

这样近似并不精确，但不这样做就无法分析摩擦力与压力的关系：决定因素很多，只能取最简单的模型。

\example{三极管}{线性近似}{}
\textbf{模型}：电路课把三极管三个引线的电流写成
\[
I,\qquad \beta I,\qquad (\beta+1)I,
\]
即看成电流按固定比例 $\beta$ 联系的线性器件。

\begin{figure}[H]
\centering
\begin{tikzpicture}[>=stealth,thick]
  \draw (0,-0.55) -- (0,0.55);
  \draw[->] (-1.5,0) -- (0,0) node[pos=0.35,above] {$I$};
  \draw[->] (0.6,1.8) -- (0.6,1.1);
  \draw (0.6,1.15) -- (0.6,0.65) -- (0,0.22);
  \node[right] at (0.7,1.35) {$\beta I$};
  \draw (0,-0.22) -- (0.6,-0.65) -- (0.6,-1.2);
  \draw[->] (0,-0.22) -- (0.45,-0.54);
  \draw[->] (0.6,-1.15) -- (0.6,-1.8);
  \node[right] at (0.7,-1.35) {$(\beta+1)I$};
  \node[red!70!black,font=\bfseries] at (3.2,0.2) {不是};
  \node[anchor=west,gray!55!black,font=\footnotesize,align=left] at (3.9,0) {实际电流并不严格\\是这样的比例};
\end{tikzpicture}
\caption{三极管的线性模型：进去 $I$，上面 $\beta I$，下面 $(\beta+1)I$。}
\end{figure}

按 NPN 三极管的标准符号，进线是\textbf{基极}(base)，上面是\textbf{集电极}(collector)，下面带箭头的是\textbf{发射极}(emitter)；$\beta$ 是\textbf{电流放大系数}(current gain)，$I_C=\beta I_B$ 正是一个线性关系。

\textbf{实际}：真实三极管并不严格满足这组线性关系，只在放大区近似如此。

不这样近似就无法分析三极管。更精确的做法是用 PSPICE、Protel 等电路软件里的\textbf{非线性}模型；手工分析或只需建立\textbf{感性认识}时，线性模型够用。

\subsection{理由二：时不变}
「时不变」的日常版：今天输入 $x(t)$ 得 $y(t)$，明天输入同样的信号 $x(t-t_0)$，得同样的输出 $y(t-t_0)$；后天、永远如此，$t_0$ 任意。

\begin{figure}[H]
\centering
\begin{tikzpicture}[>=stealth]
  \sysline{0}{今天：$x(t)$}{$y(t)$}{原信号}
  \sysline{-1.0}{明天：$x(t-t_0)$}{$y(t-t_0)$}{同样的信号，晚了 $t_0$ 才送进去}
  \sysline{-2.0}{后天、大后天……}{同样的输出}{$t_0$ 是任意实数，永远如此}
\end{tikzpicture}
\caption{时不变的日常版：同一个信号晚 $t_0$ 送进去，输出也只是晚 $t_0$ 出来，形状不变；今天、明天、永远都如此。三行里的「系统」是\textbf{同一台}机器，它的脾气不随日子变。}
\end{figure}

\textbf{绝对的时不变不存在。}若有这种系统，世界就「不老」了。有些东西能在较长时间内不变，但连宇宙都可能毁灭，所以没有绝对意义的时不变系统。

\textbf{那为什么还要假设？}三个理由：
\begin{enumerate}[nosep]
\item 不假设就无法研究：时变意味着输出随时间无规律，无法预测。
\item 研究时变系统也要先假定一个变化规律，否则无从下手。
\item 即使系统时变，在很短一段时间内也可近似为时不变（例如一块肉隔几天才腐烂，中间几小时可视为不变）。
\end{enumerate}

\begin{remember}
研究时不变系统的理由：没有绝对的时不变，但时不变是对自然的合理近似；不假设，连研究都没法开始。
\end{remember}

\subsection{理由三：LTI 很简单}
\textbf{LTI 本身很简单}，核心是下面这句：
\begin{remember}
知道 LTI 系统的\emph{一个} $x(t)$ 对应的 $y(t)$，就知道\emph{所有} $x(t)$ 对应的 $y(t)$。
\end{remember}
第二章那么多数学、公式，核心就是这一句；后面的卷积、冲激响应都在回答「怎么从那一个推出所有」。

\begin{skippable}
那「一个」不能是任意信号：取「输入恒 $0$、输出恒 $0$」，什么也推不出来。它必须是有代表性的信号——单位冲激 $\delta(t)$，其响应叫\textbf{冲激响应}。任意信号可拆成冲激的叠加（线性），每个冲激只是平移（时不变），所以知道 $\delta(t)$ 的响应就能求任意输入的输出。
\end{skippable}

%==================================================================
\needspace{6cm}
\section{与《典型的系统》的关系}
《典型的系统》讲「怎么判断一个系统有没有某个性质」；本笔记讲「为什么要关心这些性质」：给出 LTI 的定义、六个性质的英文、三个研究理由。判断方法见《典型的系统》，这里不重复。

\end{document}
